\section*{Riesgos Asociados con la Radiación}

Como ya se mencionó, la utilización de materiales radiactivos implica la presencia de alguno o varios de los riesgos siguientes:

\subsection*{Tipos de Irradiación}
\begin{itemize}
    \item \textbf{Externa:} gamma, beta
    \item \textbf{Interna:} alfa, beta, gamma
\end{itemize}

\subsection*{Contaminación}
\begin{itemize}
    \item \textbf{Externa:} gamma, beta
    \item \textbf{Interna:} alfa, beta, gamma
\end{itemize}

\section*{Liberación de Materiales Radiactivos}
Los materiales radiactivos podrán liberarse al medio ambiente cuando ocurra algún accidente que afecte la integridad del contenedor y cápsula. Cuando esto ocurre, la presencia del material radiactivo ocasiona riesgos de contaminación externa y/o interna ya que el material puede embarrarse en la piel durante las actividades de trabajo o incluso ingerirse o inhalarse, dependiendo de las características físicas del material radiactivo.

\section*{Medidas de Protección Contra la Irradiación Externa}
A continuación se mencionan algunas medidas de protección contra la irradiación externa que se deberán seguir siempre al tomar una radiografía:

\begin{itemize}
    \item Antes de empezar el trabajo, verificar que su dosímetro de lectura directa no se encuentre saturado o fuera de escala y que su dosímetro de película esté en buen estado.
    \item Verificar que el detector geiger tiene baterías, que está bien calibrado y que su alarma está encendida.
    \item Antes de empezar la irradiación se debe acordonar una área aproximada de 15 m. de la fuente con letreros que indiquen la presencia de radiación.
    \item Durante el tiempo que dure la exposición, se debe mantener alejado a la gente pero controlando el área para evitar que algunas personas puedan acercarse y recibir una exposición innecesaria.
    \item Al terminar la exposición, se debe verificar con el detector geiger que la fuente efectivamente regresó al contenedor y que se encuentra segura. Es vital tomar los niveles de radiación cerca del contenedor y de la manguera posicionadora.
    \item Volver a leer su dosímetro de bolsillo y tomar nota de la dosis que recibió durante su jornada de ese día. Al final del día, registrar su dosis en las formas ya previstas para ello y al final del año de la misma manera en registros anuales.
    \item En caso de exposiciones con grado de dificultad mayor, como por ejemplo en tuberías aéreas, en donde el técnico debe permanecer cerca de la fuente y que el tiempo de irradiación va a ser muy corto y no permite alejarse y regresar, el uso de el equipo o las indumentarias que le permitan tener que planear el trabajo en forma tal que se aplique lo que hasta aquí ha aprendido.
\end{itemize}
